\documentclass[a4paper,11pt]{article}
\usepackage[T1]{fontenc}
\usepackage[margin=2.5cm]{geometry}
\usepackage{amsmath}
\usepackage{amssymb}
\usepackage{amsthm}
\usepackage{mathtools}
\usepackage{physics}
\usepackage{bm}
\usepackage{graphicx}
\usepackage{float}
\usepackage{subcaption}
\usepackage{booktabs}
\usepackage{siunitx}
\usepackage{xcolor}
\usepackage{listings}
\usepackage[colorlinks=true,linkcolor=blue,citecolor=blue,urlcolor=blue]{hyperref}
\usepackage[nameinlink]{cleveref}

\title{This is a Title}
\author{Physics Society-Academic Officer-Hugh Rodgers}
\date{last update \today}

\begin{document}
    \maketitle
    \tableofcontents
    \newpage
    \raggedright

    This document comes with a .TEX file which contains all the LaTeX syntax, it is intended to be read at the same time as the PDF.
    Comments in latex are done as a  percent sign (\%) followed by the comment text. \\ 
    \% Insert comment here\\
    The accompanying .TEX file has many comments explaining the LaTeX syntax used.
    % This is what a comment looks like, you may notice it does not appear in the PDF

    \section{Basic Text}
        This is normal text.\\
        \textbf{Bold text}\\
        \textit{Italic text}\\
        \underline{Underlined text}
        \begin{center}
            Centered Text
        \end{center}

        \begin{flushright}
            Right-aligned Text
        \end{flushright}

        \begin{flushleft}
            Left-aligned Text
        \end{flushleft}

        \Huge{This is big text}\\
        \LARGE{This is larger text}\\
        \Large{This is large text}\\
        \normalsize{This is normal text}\\
        \small{This is small text}\\
        \footnotesize{This is footnote size text}\\
        \scriptsize{This is script size text}\\
        \tiny{This is tiny text}\\

        \normalsize
        If you want two columns, you can use the \texttt{\textbackslash twocolumn[]} and return to single column with \texttt{\textbackslash onecolumn}.
        % the \tttext is used to display text in typewriter (monospaced) font and \backslash is used to display the backslash character itself on the pdf
        % the backslash in latex usually indicates a command meaning what follows is given special meaning

    \section{Lists} % it is generaly not a good idea to use this in your labreport unless completely necessary, 
        \begin{itemize}
        \item First item
        \item Second item
        \end{itemize}

        \begin{enumerate}
        \item Step One
        \item Step Two
        \end{enumerate}


    \section{Equations}
        There are many different ways of writing equations in LaTeX.

        \text{\$ F=ma\$} returns - $F=ma$ - which is in the inline math mode meaning writing equations within a line of text.

        The rest of the equations are hard to write in the pdf so please look at the .tex file to see how they are written.
        \begin{equation}
        E = mc^2
        \label{eq:newton}
        \end{equation}
        This makes the equation referenceable in the document~\ref{eq:newton}.

        \[
        \left\langle \psi_i |\psi_j \right\rangle = \delta_{i,j}
        \]

        \(
        u^i v_i = \sum_i^n u^i v_i = \mathbf{u} \cdot \mathbf{v} = u_1 v_1 + u_2 v_2 + 
        \dots + u_n v_n
        \)
        
        \[
        \dv{y}{x},\qquad % \qquad is a method to add space between elements
        \pdv{f}{x},\qquad % other methods spacing between elements in math mode are \, \,, \:, \;, \quad, \qquad
        \int_0^1 x^2\,dx,\qquad
        \iiint_0^1 x^2\,dxdydz
        \]

        \[
        A=\begin{pmatrix}
        1 & 2\\
        3 & 4
        \end{pmatrix}
        \]

        \begin{table}[H]
        \centering
        \begin{tabular}{cc}
        \toprule
        x & y\\
        \midrule
        1 & 2\\
        2 & 4\\
        \bottomrule
        \end{tabular}
        \caption{Example Table}
        \end{table}

        \[
        c=\SI{2.99792458e8}{\metre\per\second}
        \]

        \[
        \vb{F}=m\vb{a}
        \]

        \[
        \nabla\cdot\vb{E}=\frac{\rho}{\epsilon_0}
        \]

 
    \section{Code Example} % code should fall in the appendex of the lab report just to show what was done, may not be necessary but easier that a screenshot 
        \begin{lstlisting}
        import numpy as np
        x = np.linspace(0,10,100)
        print(x)
        \end{lstlisting}


    \section{Figures}
        Add figures using:

        % [H] forces the figure to appear exactly here
        \begin{figure}[H]
            % other things can be included to make you image look nicer
            \centering
            \includegraphics[width=0.5\textwidth]{LaTeX_meme.png}
            \caption{Example figure Caption}
            \label{fig:example_figure_label}

        \end{figure}
        \ref{fig:example_figure_label}

    \section{Links}
        This is an example of a link to Overleaf LaTeX Examples \cite{overleaf_examples}


    \section{references and bib files}

        \bibliographystyle{plain}
        \bibliography{LaTeX/ref.bib}

\end{document}
